% !PRINTED 161
\weberpage{161}
%<<CONT
über, woraus man erkennt, dass $\alpha$ und $\gamma'$ die beiden dritten Wurzeln $\varepsilon$ und $\varepsilon^2$ der Einheit sind. Die vier gesuchten Formen $f$ sind demnach:
\[
\begin{aligned}
f_0 &= 2\,x_1x_2\\
f_1 &= x_1^2 - \frac{2}{3}\,x_1x_2 - \frac{8}{9}\,x_2^2\\
f_2 &= \varepsilon\,x_1^2 - \frac{2}{3}\,x_1x_2 - \frac{8}{9}\,\varepsilon^2 x_2^2\\
f_3 &= \varepsilon^2\,x_1^2 - \frac{2}{3}\,x_1x_2 - \frac{8}{9}\,\varepsilon\,x_2^2\,.
\end{aligned}
\]
Damit ist die Aufgabe gelöst, quadratische Formen $f$ so zu bestimmen, dass ihre zweiten Ueberschiebungen $A_{ik}$ durchwegs den Werth $+\varrho$ besitzen.

\textit{Anmerkung.} Das Product der vier quadratischen Formen liefert eine Form achten Grades:
\[
W = \text{Const.}\ xy\,(x^6 - 8x^3y^3 - 8y^6)\,,
\]
wenn wir der Einfachheit halber $\frac{2}{3}x_2=y$, $x_1=x$ setzen. Klein nennt diese Form (vgl. Vorles. über das Ikosaeder \S~5 und 10) Würfel, da man wiederum die Gauss'sche Zahlenebene so auf eine im Nullpunkt derselben berührende Kugel projiciren kann, dass die acht Wurzeln von $W=0$ auf der Kugel acht Punkte bestimmen, die gerade die acht Ecken eines der Kugel einbeschriebenen Würfels bilden. Ueber diese Form $W$ vergleiche auch \S~19.

142. \textit{Relationen zwischen den vier Formen $f_i$.} Man erkennt direct, dass die Summe aller vier Formen identisch verschwindet; demnach hat man als erste Relation:
\weq{1}{f_0+f_1+f_2+f_3=0\,.}
Aus ihr folgt, dass die Summe irgend zweier derselben bis auf das Vorzeichen gleich der Summe der beiden andern ist. Bezeichnen wir eine solche Summe mit $\varphi_{ik}$, so hat man:
\weq{2}{\begin{aligned}
\varphi_{01} &= -\varphi_{23}\\
\varphi_{02} &= -\varphi_{13}\\
\varphi_{03} &= -\varphi_{12}\,.
\end{aligned}}
Da also die sechs Summen zu je zweien der vier Formen $f_i$ nur drei verschiedene absolute Werthe annehmen, so müssen sie sich als Wurzeln einer Gleichung sechsten Grades darstellen lassen, die sich durch die einfache Substitution
\[
\varphi_{ik}^2 = z
\]
auf eine' cubische $\Phi=0$ reduciren lässt. % [sic] stray mark printed after \textquotedbl{}eine\textquotedbl{}


% ===== pdf idx 176, printed 162.0 =====
% !PRINTED 162
\weberpage{162}
Die drei Formen $\varphi_{01}$, $\varphi_{02}$, $\varphi_{03}$ sind drei conjugirte quadratische Formen; denn man hat:
\weq{3}{\begin{aligned}
(\varphi_{01},\varphi_{02})^2 &= 0\\
(\varphi_{02},\varphi_{03})^2 &= 0\\
(\varphi_{03},\varphi_{01})^2 &= 0\,.
\end{aligned}}
Dies ergiebt sich sofort durch Rechnung. So ist beispielsweise:
\[
\begin{aligned}
(\varphi_{01},\varphi_{02})^2 &= (f_0+f_1,\,f_0+f_2)^2 = A_{00}+A_{01}+A_{02}+A_{12}\\
&= -2 + \frac{2}{3} + \frac{2}{3} + \frac{2}{3} = 0\,.
\end{aligned}
\]
Es ist ferner bemerkenswerth, dass auch eine lineare Relation zwischen den Quadraten der vier Formen $f_i$ existirt. Man verificirt nämlich sehr leicht durch Rechnung:
\weq{4}{f_0^2+f_1^2+f_2^2+f_3^2=0\,.}

\textit{Anmerkung.} 1. Wegen der Relationen (1) und (4) reducirt sich die Gleichung:
\[
(\xi-f_0)(\xi-f_1)(\xi-f_2)(\xi-f_3)=0
\]
auf die drei Glieder:
\[
\xi^4+A\xi+B=0\,.
\]
Durch die Substitution
\[
\xi=\eta\sqrt[3]{A}
\]
geht sie über in die Gleichung
\weq{I}{\eta^4+\eta+c=0\,,}
eine Gleichung vierten Grades, die nur einen einzigen Parameter $C$ besitzt. Sie hat eine cubische Resolvente, die oben erwähnte Gleichung $\Phi=0$, die natürlich ebenfalls nur diesen einen Parameter $C$ enthalten kann. Man könnte diesen Umstand benutzen — wenn man nicht aus praktischen Gründen davon absehen müsste — die allgemeine Gleichung vierten Grades aufzulösen. Zu dem Zwecke bringt man sie zunächst — was immer möglich ist — auf die Form (I), vergleicht die Parameter und löst nun die Gleichung (I) mit Hilfe der cubischen Resolvente, aus deren Wurzeln $\varphi_{ik}$ sich leicht die Wurzeln $f_i$ der vorgelegten Gleichung bestimmen lassen. (Vgl. auch \S~19.)

\textit{Anmerkung.} 2. Herr Brill hat im 20. Bande der Math. Annalen (Binäre Formen und Gleichung sechsten Grades) gleichfalls gezeigt: Man kann immer vier quadratische Formen so finden, dass zwischen ihnen sowohl als zwischen ihren Quadraten je eine lineare Relation besteht, also dass man, wenn $f_1$ $f_2$ $f_3$ $f_4$ die vier quadratischen Formen, die Beziehungen hat:
\weq{1}{f_1+f_2+f_3+f_4=0}
\weq{2}{\alpha_1f_1^2+\alpha_2f_2^2+\alpha_3f_3^2+\alpha_4f_4^2=0\,.}


% ===== pdf idx 177, printed 163.0 =====
% !PRINTED 163
\weberpage{163}
Er zeigt ferner, dass, wenn man irgend drei der vier Formen durch ihre drei Functionaldeterminanten $\vartheta_{ik}$ ersetzt, zwischen dieser und der vierten Form ebenfalls analoge Relationen bestehen. Da dies auf vier verschiedene Arten möglich ist, so haben folgende fünf Formenquadrupel
\[
\begin{array}{cccc}
f_1 & f_2 & f_3 & f_4\\
f_1 & \vartheta_{23} & \vartheta_{34} & \vartheta_{42}\\
\vartheta_{13} & f_2 & \vartheta_{34} & \vartheta_{41}\\
\vartheta_{12} & \vartheta_{24} & f_3 & \vartheta_{41}\\
\vartheta_{12} & \vartheta_{23} & \vartheta_{31} & f_4
\end{array}
\]
die nämliche Eigenschaft, gleichzeitig durch Relationen von der Form (1), (2) aneinander gebunden zu sein. Ist also ein Quadrupel gefunden, so erhält man aus ihm stets vier andere. Hiebei kann man jedes der fünf Quadrupel in irgend einer Permutation zu Grunde legen, um auf dieselben vier andern zu gelangen. Man hat also $24\cdot5=120$ Operationen zur Verfügung, um von einem Quadrupel auf ein anderes der fünf zu gelangen, und insofern bilden die fünf Quadrupel eine Gruppe im Sinne der Gleichungstheorie.

142. % [sic] section number \textquotedbl{}142.\textquotedbl{} repeats the one used earlier in this fragment for \textquotedbl{}Relationen zwischen den vier Formen f_i\textquotedbl{} — reproduced as printed
\textit{Die quadratischen Formen des Ikosaeders.} Indem wir der Invariante $A_{f_2f_3}$ die Bedingung auferlegten, dass auch sie in Uebereinstimmung mit den fünf übrigen den constanten Werth $+\varrho$ besitzen möge, gelangten wir zu vier quadratischen Formen, die durch diese Bedingung völlig bestimmt waren. Es stand uns damals aber frei, für $A_{f_2f_3}$ den Werth $\pm\varrho$ zu wählen, und wir wollen nun untersuchen, zu welchem Resultate die Wahl $A_{f_2f_3}=-\varrho$ führt. In diesem Falle wird die Bedingung für den numerischen Werth von $\varrho$ dargestellt durch die Determinante
\[
\begin{vmatrix}
-2 & \varrho & \varrho & \varrho\\
\varrho & -2 & \varrho & \varrho\\
\varrho & \varrho & -2 & -\varrho\\
\varrho & \varrho & -\varrho & -2
\end{vmatrix} = 0\,.
\]
Wir ersetzen in ihr wiederum $-\dfrac{2}{\varrho}$ durch $\sigma$ und erhalten alsdann
\[
\begin{aligned}
\begin{vmatrix}
\sigma & 1 & 1 & 1\\
1 & \sigma & 1 & 1\\
1 & 1 & \sigma & -1\\
1 & 1 & -1 & \sigma
\end{vmatrix}
&=
\begin{vmatrix}
\sigma-1, & 1-\sigma, & 0, & 0\\
1, & \sigma, & 1, & 1\\
0, & 0, & (\sigma+1), & -(\sigma+1)\\
1, & 1, & -1, & \sigma
\end{vmatrix}\\[4pt]
&= (\sigma^2-1)
\begin{vmatrix}
1 & -1 & 0 & 0\\
1 & \sigma & 1 & 1\\
0 & 0 & 1 & -1\\
1 & 1 & -1 & \sigma
\end{vmatrix} = 0\,.
\end{aligned}
\]


% ===== pdf idx 178, printed 164.0 =====
% !PRINTED 164
\weberpage{164}
Der Factor $\sigma^2-1=0$ führt wiederum auf den Werth $\varrho=\pm 2$, den wir aus gleichen Gründen wie früher verwerfen müssen. Der andere Factor reducirt sich auf
\[
(\sigma+1)(\sigma-1)-4=0
\]
oder
\[
\sigma=\pm\sqrt5\,.
\]
Die Wahl des Vorzeichens steht uns frei. Wir nehmen zunächst
\[
\sigma=+\sqrt5
\]
und erhalten alsdann
\[
\varrho=\frac{-2}{\sqrt5}\,.
\]
Die Coefficientenrelationen (I) in Nr. 140 gehen dann, wenn wir diesmal den noch willkürlichen Coefficienten von $x^2$ in $f_1$ gleich $\dfrac{2}{\sqrt5}$ wählen, anstatt ihn wie früher gleich 1 zu setzen, über in:
\weq{III}{\left.\begin{aligned}
\frac{2}{\sqrt5}\,c - b^2 &= -1\\
b &= \frac{1}{\sqrt5}\\
\alpha\gamma-\beta^2 &= -1\\
\beta &= \frac{1}{\sqrt5}\\
\frac{2}{\sqrt5}\,\gamma+c\alpha-2b\beta &= \frac{-2}{\sqrt5}
\end{aligned}\right\}}
Hieraus folgt zunächst:
\[
b=\frac{1}{\sqrt5}\,,\qquad \beta=\frac{1}{\sqrt5}\,,\qquad c=\frac{-2}{\sqrt5}\,.
\]
Die übrigen Gleichungen gehen durch Substitution dieser Werthe über in:
\[
\frac{2}{\sqrt5}\,\gamma-\frac{2}{\sqrt5}\,\alpha-\frac{2}{5}=\frac{-2}{\sqrt5}
\]
\[
\alpha\gamma=\frac{-4}{5}\,.
\]
Wenn wir hier wiederum, um die Gleichungen symmetrischer zu gestalten, in $f_2$ den Coefficienten $\alpha$ durch $\dfrac{2}{\sqrt5}\alpha'$ und den Coefficienten $\gamma$ durch $\dfrac{-2}{\sqrt5}\gamma'$ ersetzen, so erhalten wir:
\weq{IV}{\left.\begin{aligned}
\alpha'+\gamma' &= \frac{-1+\sqrt5}{2}\\
\alpha'\gamma' &= 1
\end{aligned}\right\}}


% ===== pdf idx 179, printed 165.0 =====
% !PRINTED 165
\weberpage{165}
Nun sind bekanntlich die fünften Wurzeln der Einheit
\[
\begin{aligned}
\varepsilon &= \dfrac{-1+\sqrt5}{4} + i\,\dfrac{\sqrt{10+2\sqrt5}}{4}, &
\varepsilon^3 &= \dfrac{-1-\sqrt5}{4} - i\,\dfrac{\sqrt{10-2\sqrt5}}{4} \\
\varepsilon^2 &= \dfrac{-1-\sqrt5}{4} + i\,\dfrac{\sqrt{10-2\sqrt5}}{4}, &
\varepsilon^4 &= \dfrac{-1+\sqrt5}{4} - i\,\dfrac{\sqrt{10+2\sqrt5}}{4} \\
\varepsilon^5 &= 1. &&
\end{aligned}
\]
Demnach kann man die Gleichungen (IV) auch in der Form schreiben
\weq{V}{\begin{aligned} \alpha' + \gamma' &= \varepsilon + \varepsilon^4 \\ \alpha'\,\gamma' &= \varepsilon\cdot\varepsilon^4 \end{aligned}}
oder, wenn man auch in den Gleichungen (IV) der $\sqrt5$ das negative Zeichen giebt, was nach unseren Entwicklungen erlaubt ist,
\weq{VI}{\begin{aligned} \alpha' + \gamma' &= \varepsilon^2 + \varepsilon^3 \\ \alpha'\,\gamma' &= \varepsilon^2\cdot\varepsilon^3. \end{aligned}}
Die beiden Gleichungen (V) und (VI) aber lehren:

\glqq Die Coefficienten $\alpha'$ und $\gamma'$ der quadratischen Form $f_2$ sind je ein Paar $\varepsilon^\nu$ und $\varepsilon^\mu$ der fünften Einheitswurzeln, wenn man dieselben so auswählt, dass die Exponentensumme $\nu+\mu=5$ ist.\grqq{}

Es existiren demnach vier Formen $f_2$, welche die verlangten Eigenschaften besitzen. Wir erhalten also sechs quadratische Formen von der Beschaffenheit, dass ihre fünfzehn zweiten Ueberschiebungen $A_{f_if_k}$ $(i\gtrless k)$ bis auf das Vorzeichen übereinstimmen. Sie sind, wenn wir die Formen $f_1,\dots f_5$ noch mit $\dfrac{\sqrt5}{2}$ multipliciren, dargestellt durch:
\[
\begin{aligned}
f_0 &= 2x_1x_2 \\
\frac{\sqrt5}{2}f_1 &= x_1^2 + x_1x_2 - x_2^2 \\
\frac{\sqrt5}{2}f_2 &= \varepsilon\,x_1^2 + x_1x_2 - \varepsilon^4 x_2^2 \\
\frac{\sqrt5}{2}f_3 &= \varepsilon^2 x_1^2 + x_1x_2 - \varepsilon^3 x_2^2 \\
\frac{\sqrt5}{2}f_4 &= \varepsilon^3 x_1^2 + x_1x_2 - \varepsilon^2 x_2^2 \\
\frac{\sqrt5}{2}f_5 &= \varepsilon^4 x_1^2 + x_1x_2 - \varepsilon\, x_2^2.
\end{aligned}
\]
\textit{Anmerkung.} Das Product dieser sechs quadratischen Formen bezeichnete Klein mit dem Namen \glqq Ikosaeder\grqq{} (vergl. a. a. O.), aus den nämlichen Gründen, die ihn zu den Bezeichnungen \glqq Oktaeder\grqq{} und \glqq Würfel\grqq{} veranlassten. Die Gleichung des Ikosaeders ist demnach
%CONT>>


% ===== pdf idx 180, printed 166.0 =====
% !PRINTED 166
\weberpage{166}
%<<CONT
\[
x_1x_2\big(x^5 - (\varepsilon+\varepsilon^4)^5 x_2^5\big)\big(x^5 - (\varepsilon^2+\varepsilon^3)^5 x_2^5\big) = 0\,, % [sic] printed \textquotedbl{}x^5\textquotedbl{} without subscript in the first two factors
\]
oder:
\[
x_1x_2\big(x_1^{10} + 11\,x_1x_2 - x_2^{10}\big) = 0.
\]
Ich möchte noch darauf aufmerksam machen, dass auch zwischen den sechs quadratischen Formen specielle Relationen bestehen müssen, und zwar vier an der Zahl, da sie nur von den zwei Parametern $x_1$ und $x_2$ abhängen. In der That findet man leicht durch directe Rechnung, wenn man die erste Form $f_0$ noch mit $\dfrac{\sqrt5}{2}$ multiplicirt und nun $\dfrac{\sqrt5}{2}f_i$ mit $\psi_i$ bezeichnet, die vier Identitäten:
\weq{1}{\sum_{i=0}^{i=5}\psi_i^2 = 0\,,}
\weq{2}{\sum_{i=0}^{i=5}\psi_i^4 = 0\,,}
\weq{3}{\sum_{i=0}^{i=5}\psi_i^6 = 0\,,}
\weq{4}{\sum_{i=0}^{i=5}\psi_i^3 = \dfrac{(1+\sqrt5)^3}{(1-\sqrt5)^2\sqrt5}\Big(\sum_1^5 \psi_i\Big)^3.}
Aus den ersten drei Relationen geht sofort hervor, dass die Gleichung
\[
\prod_{i=0}^{i=5}(z-\psi_i^2) = 0
\]
die Form haben wird
\[
z^6 + az^2 + bz + c = 0.
\]
Durch die Transformation $z = \sqrt[4]{a}\cdot y$ geht sie über in die Form
\weq{5}{y^6 + y^2 + b'y + c' = 0.}
Aber auch diese beiden Constanten $b'$ und $c'$ sind noch von einander abhängig wegen der noch unbenutzten Relation (4). Sie sind beide Functionen ein und desselben Parameters $\lambda$, so dass also (5) eine Gleichung sechsten Grades mit einem \gesperrt{einzigen} Parameter darstellt. Wir werden später auf solche Gleichungen noch zu sprechen kommen. Nimmt man statt der sechs quadratischen Formen ihre ersten Polaren $f_{i_y}$, so bestehen immer noch vier analoge Identitäten, wie die oben angeführten. Aber die Gleichung sechsten Grades, welche die Quadrate dieser Formen zu Wurzeln hat, besitzt nunmehr zwei Parameter. Es ist dies dieselbe Gleichung sechsten Grades, welche Jacobi zwischen Thetareihen im 2. Bd. des Crelle'schen Journals erhalten hat. Setzt man nämlich:
%CONT>>


% ===== pdf idx 181, printed 167.0 =====
% !PRINTED 167
\weberpage{167}
%<<CONT
\[
\Theta(q) = 1 + 2q + 2q^4 + 2q^9 + 2q^{16} + 2q^{25} + \cdots
\]
und
\[
\varepsilon = e^{\frac{2i\pi}{5}}\,,
\]
so wird
\[
\Theta(\varepsilon^\varrho\cdot q) = 1 + 2\varepsilon^\varrho q + 2\varepsilon^{4\varrho}q^4 + 2\varepsilon^{4\varrho}q^9 + 2\varepsilon^\varrho q^{16} + 2q^{25} + \cdots\,,
\]
oder
\[
\Theta(\varepsilon^\varrho q) = \varepsilon^\varrho A_0 + A_1 + \varepsilon^{4\varrho}A_2\,,
\]
wobei
\[
\begin{aligned}
A_0 &= 2q + 2q^{16} + 2q^{36} + 2q^{81} + \cdots \\
A_1 &= 1 + 2q^{25} + 2q^{100} + 2q^{225} + \cdots \\
A_2 &= 2q^4 + 2q^9 + 2q^{49} + 2q^{64} + \cdots\,.
\end{aligned}
\]
Setzen wir aber $A_0 = x_1y_1$, $A_1 = (x_1y_2 + x_2y_1)$, $A_2 = -x_2y_2$, so haben wir:
\[
\begin{aligned}
\Theta(q) &= f_{1y} \\
\Theta(\varepsilon q) &= f_{2y} \\
\Theta(\varepsilon^2 q) &= f_{3y} \\
\Theta(\varepsilon^3 q) &= f_{4y} \\
\Theta(\varepsilon^4 q) &= f_{5y}\,,
\end{aligned}
\]
und überdies:
\[
\Theta(q^{25}) = f_{0y}\,.
\]
Es lassen sich also in der That diese sechs Thetareihen mit den Polaren der sechs quadratischen Formen $f_0, f_1, \dots f_5$ identificiren.

\begin{center}
\S~14.\\[2pt]
\textbf{Die Form dritten Grades.}
\end{center}

\textbf{143.} \textit{Einleitende Bemerkungen.} Wir haben bereits in den Beispielen Nr. 37 und Nr. 39, welche den Ueberschiebungsprocess erläutern sollten, eine Reihe von Co- und Invarianten der Form dritten Grades kennen gelernt. Bezeichnet man die cubische Form mit
\[
\begin{aligned}
f &= \bar a_0 x_1^3 + 3\bar a_1 x_1x_2 + 3\bar a_2 x_1x_2^2 + \bar a_3 x_2^3 \\
&= a_x^3 = b_x^3 = c_x^5 = \text{etc.}\,,
\end{aligned}
\]
% [sic] printed exponent "c_x^5"; the other symbolic powers in the line read a_x^3 = b_x^3
so haben die a. a. O. gebildeten Covarianten die symbolischen Producte
\[
\begin{aligned}
(f,f)^2 &= (ab)^2\,a_xb_x = \Delta \\
(\Delta,\Delta)^2 &= (ab)^2(cd)^2(ad)(bc) = A_{\Delta\Delta} = R \\
(f,\Delta) &= (ab)^2(ca)\,b_x c_x^2 = Q
\end{aligned}
\]
zum Repräsentanten. Schon damals bemerkten wir, dass mit diesen vier Formen
\[
f,\ \Delta,\ R,\ Q
\]
das volle System von $f$ abgeschlossen ist, d. h. dass alle übrigen Co- und Invarianten von $f$ sich rational und ganz durch dieselben dar-
%CONT>>


% ===== pdf idx 182, printed 168.0 =====
% !PRINTED 168
\weberpage{168}
%<<CONT
stellen lassen, und es ist unsere erste Aufgabe, diese Behauptung nachzuweisen. Da nun alle weiteren Co- und Invarianten nur durch Ueberschiebung der drei Covarianten $f$, $\Delta$, $Q$ entstehen können und somit durch symbolische Producte dargestellt sein müssen, die irgend welche der sechs Klammerfactoren
\[
(ab),\ (a\Delta),\ (\Delta\Delta_1),\ (aQ),\ (\Delta Q),\ (QQ_1)
\]
zu Factoren haben, so besteht unsere Aufgabe darin, den Nachweis zu liefern, dass jeder dieser sechs Klammerfactoren \gesperrt{Reducent} ist (vergl. Nr. 124), d. h. dass jedes symbolische Product, welches solche Klammerfactoren besitzt, entweder identisch verschwindet oder auf Producte mit Factoren $f$, $\Delta$, $R$, $Q$ reducirt werden kann. Indem ich nun diesen Beweis antrete, will ich nur den Beweis für die Reducenteneigenschaft des Klammerfactors $(ab)$ nochmals in voller Ausführlichkeit entwickeln. Da der Gedankengang bei den folgenden Beweisen derselbe bleibt, so wird alsdann eine grössere Kürze gestattet sein.

\textbf{144.} \textit{Der Factor $(ab)$ ist Reducent.} Jedes symbolische Product $P$, das den Factor $(ab)$ besitzt, muss in irgend welchen Verbindungen auch noch zwei weitere Symbole $a$ und zwei Symbole $b$ enthalten, wenn es eine wirkliche Covariante von $f$ darstellen soll. Wir drücken dies am allgemeinsten aus, wenn wir sagen:
\[
\pi = (ab)\,a_\xi a_x b_\eta b_y
\]
muss ein Factor des zu untersuchenden Productes $P$ sein, wobei $\xi$, $x$, $\eta$, $y$ durch irgend welche Symbole ersetzt sein können. Dieses Theilproduct $\pi$ kann als eine gemischte Polare der Form
\[
\Pi = (ab)\,a_x^2 b_y^2
\]
angesehen werden und ist mit dieser gleichzeitig reducibel. Die Form $\Pi$ ist aber ein Glied der zweiten Polare von
\[
(f,f) = (ab)\,a_x^2 b_x^2
\]
und lässt sich somit darstellen durch diese Polare $(f,f)_{y^2}$ selbst plus Gliedern, welche den Factor $(ab)^2(xy)$ besitzen (vergl. Nr. 25). $\Pi$ ist also reducibel, sobald $(f,f)_{y^2}$ und \gesperrt{auch} die Glieder mit dem Factor $(ab)^2(xy)$, oder was dasselbe ist, sobald $(f,f)$ und $(ab)^2a_xb_y$ reducibel sind. Der Term $(ab)^2a_xb_y$ ist jedoch ein Glied der ersten Polare von
\[
(f,f)^2 = (ab)^2a_xb_x
\]
und lässt sich somit ersetzen durch diese plus Gliedern mit dem Factor $(ab)^3(xy)$. Die Reducibilität von $\Pi$ ist also zurückgeführt auf die drei Bedingungen, dass $(f,f)$, $(f,f)^2$, $(f,f)^3$ reducibel sind, d. h.
%CONT>>


% ===== pdf idx 183, printed 169.0 =====
% !PRINTED 169
\weberpage{169}
%<<CONT
auf bekannte Formen führen. Es ist aber
\[
(f,\,f) = 0, \quad (f,\,f)^2 = \Delta, \quad (f,\,f)^3 = 0,
\]
und demnach ist $\pi$ reducibel, und $(ab)$ Reducent.
% checked against the GDZ scan (27 Oct 2026): the print has a lowercase \pi here (the product \Pi of p. 168, misprinted), not \Delta; reproduced as printed

\textit{Anmerkung.} Ich möchte noch bemerken, dass sich diese Darstellung der Reducibilität mit der in Nr. 124 gegebenen vollständig deckt. Wir belegten damals Klammerfactoren mit dem Namen Reducenten, sobald das Product, welches sie allein besass, sammt allen daraus durch Faltung entstehenden reducibel war. In der That sind wir hier zum gleichen Resultate gelangt, denn $(f,\,f)^2$ und $(f,\,f)^{3}$ entstehen aus $(f,\,f)$ durch Faltung.

145. \textit{Die Factoren $(a\Delta)$ und $(\Delta\Delta_1)$ sind Reducenten.} Die erste Ueberschiebung von $f$ über $\Delta$ lieferte die Covariante $Q$; die zweite Ueberschiebung aber verschwindet identisch (vergl. auch Nr. 62, Beispiele). Denn ersetzt man in
\[
\Delta_x^2 = (ab)^2 a_x b_x
\]
$x$ durch $c$, so kommt:
\[
(c\Delta)^2 = (ab)^2 (ac)(bc),
\]
und somit:
\[
(f,\,\Delta)^2 = (c\Delta)^2 c_x = (ab)^2 (ac)(bc) c_x.
\]
Vertauscht man rechts einmal $c$ mit $b$, dann auch $c$ mit $a$ und addirt die drei so erhaltenen Ausdrücke für $(f,\,\Delta)^2$, so ergiebt sich:
\[
(f,\,\Delta)^2 = \frac{1}{3}(ab)(ac)(bc)\{(ab)c_x - (ac)b_x + (bc)a_x\},
\]
oder:
\weq{I}{(f,\,\Delta)^2 = 0,}
da der in der Klammer befindliche Ausdruck identisch verschwindet. Weitere Ueberschiebungen von $f$ über $\Delta$ existiren nicht; $(f,\,\Delta) = Q$, $(f,\,\Delta)^2 = 0$, folglich ist $(a\Delta)$ Reducent.

Ebenso ist $(\Delta\Delta_1)$ Reducent, weil $(\Delta,\,\Delta) = 0$, $(\Delta,\,\Delta)^2 = A_{\Delta\Delta}$ ist.

146. \textit{Der Factor $(aQ)$ ist Reducent.} Nach den Entwicklungen in Nr. 144 haben wir nur den Werth der drei Ueberschiebungen:
\[
(f,\,Q), \quad (f,\,Q)^2 \quad \text{und} \quad (f,\,Q)^3
\]
zu untersuchen. Wir schlagen hiebei ganz den regulären Gang des Ueberschiebungsprocesses ein, indem wir zunächst die erste Polare von $Q$ bilden. Es wird:
\weq{1}{3\,Q_x^2 Q_y = (c\Delta)\{c_x^2 \Delta_y + 2\Delta_x c_x c_y\} = G_1 + 2G_2.}
Wir können den Ausdruck rechts vereinfachen. Die Differenz aus der Polare und einem ihrer Glieder, nämlich
\[
3\,Q_x^2 Q_y - 3G_2 = (c\Delta)\{c_x^2 \Delta_y + 2\Delta_x c_x c_y\} - 3(c\Delta)\Delta_x c_x c_y,
\]
%CONT>>


% ===== pdf idx 184, printed 170.0 =====
% !PRINTED 170
\weberpage{170}
%<<CONT
hat den Factor $(c\Delta)^2(xy)$. Ein Product mit dem Factor $(c\Delta)^2$ ist aber identisch Null, da $(f,\,\Delta)^2 = 0$, und ein solches Product aus dieser Form durch Ueberschiebung entstanden gedacht werden kann. Folglich sind die beiden Glieder in (1) einander gleich, d. h. es ist:
\[
(c\Delta)c_x^2 \Delta_y = (c\Delta)\Delta_x c_y c_x
\]
und demnach wird die Polare:
\weq{2}{Q_x^2 Q_y = (c\Delta)c_x^2 \Delta_y = (c\Delta)c_x c_y \Delta_x.}
Hieraus folgt für $y_1 = -b_2$, $y_2 = b_1$:
\[
(bQ)Q_x^2 b_x^2 = (c\Delta)(bc)c_x \Delta_x b_x^2,
\]
oder, indem man $b$ mit $c$ vertauscht und beide Ausdrücke addirt:
\[
(f,\,Q) = \frac{1}{2}(bc)c_x b_x \Delta_x \big((c\Delta)b_x - (b\Delta)c_x\big)
= -\frac{1}{2}(bc)^2 b_x c_x \Delta_x^2,
\]
oder
\weq{I}{(f,\,Q) = -\frac{1}{2}\Delta^2.}
Bildet man ferner die erste Polare von $\Delta_x$, so kommt:
\weq{3}{\Delta_y \Delta_x = (ab)^2 a_x b_y.}
Wir ersetzen hierin $x$ durch $\Delta'$ und erhalten nach Multiplication mit $\Delta'_y$:
\[
(\Delta\Delta')\Delta_y \Delta'_y = (ab)^2(a\Delta')b_y \Delta'_y = (ab)^2(a\Delta)b_y \Delta_y.
\]
Substituiren wir dagegen in (2) für $x$ das Symbol $b$, so kommt:
\[
(Qb)^2 Q_y b_y = (cb)^2(c\Delta)b_y \Delta_y.
\]
Durch Comparation dieser beiden Gleichungen ergiebt sich:
\[
(Qb)^2 Q_y b_y = (\Delta\Delta')\Delta_y \Delta'_y.
\]
Rechts steht die erste Ueberschiebung von $\Delta_y^2$ über sich selbst, sie verschwindet identisch; also hat man:
\weq{II}{(f,\,Q)^2 = 0.}
Setzt man endlich in $Q = (c\Delta)c_x^2 \Delta_x$ das Symbol $b$ an die Stelle von $x$, so erhält man:
\[
(Qb)^3 = (c\Delta)(cb)^2(\Delta b).
\]
Dasselbe ergiebt sich, wenn man in $\Delta = (cb)^2 c_x b_x$ an Stelle von $x$ das Symbol $\Delta$ substituirt, nämlich:
\[
(\Delta,\,\Delta)^2 = (c\Delta)(cb)^2(b\Delta).
\]
Durch Comparation ergiebt sich also:
\weq{III}{(f,\,Q)^3 = -(Qb)^3 = (\Delta,\,\Delta)^2 = A_{\Delta\Delta}.}
Die drei Ueberschiebungen von $f$ über $Q$ sind somit auf bekannte Formen reducibel und demnach ist $(aQ)$ Reducent.


% ===== pdf idx 185, printed 171.0 =====
% !PRINTED 171
\weberpage{171}
147. \textit{Die Factoren $(\Delta Q)$ und $(QQ')$ sind Reducenten.} Der Nachweis, dass auch diese letzten Factoren Reducenten sind, lässt sich nun auf Grund der eben gewonnenen Beziehungen einfach führen. Da nach Nr. 146 (2):
\weq{1}{Q_x^2 Q_y = (c\Delta)c_x^2 \Delta_y,}
so hat man hieraus für $y_1 = -\Delta_2$, $y_2 = +\Delta_1$
\[
(\Delta Q)Q_x^2 \Delta_x = (c\Delta)(\Delta'\Delta)c_x^2 \Delta'_x
= \frac{1}{2}(\Delta'\Delta)c_x^2 \{(c\Delta)\Delta'_x - (c\Delta')\Delta_x\}
= +\frac{1}{2}(\Delta'\Delta)^2 \cdot c_x^3,
\]
oder
\weq{IV}{(\Delta,\,Q) = \frac{1}{2}\cdot A_{\Delta\Delta}\cdot f.}
Setzt man in derselben Relation (1) statt $x$ das Symbol $\Delta (= \Delta')$, so kommt:
\[
(Q\Delta)^2 Q_y = (c\Delta)(c\Delta')^2 \Delta_y = 0,
\]
da rechts der Reducent $(c\Delta)^2$ auftritt (vergl. Nr. 145 und 146). Demnach ist:
\weq{V}{(\Delta,\,Q)^2 = 0}
und folglich ist $(\Delta Q)$ Reducent.

Substituirt man endlich $Q'$ in diese Relation (1) für die Variable $x$, und multiplicirt mit $Q'_y (= Q_y)$, so erhält man:
\[
(QQ')^2 Q_y Q'_y = (c\Delta)(cQ)^2 \Delta_y Q_y.
\]
Der symbolische Ausdruck rechts besitzt aber den Werth $\frac{1}{2}A_{\Delta\Delta}\cdot\Delta$. Denn die erste Polare der verschwindenden Ueberschiebung $(f,\,Q)^2 = (aQ)^2 a_x Q_x$ ist:
\[
(f,\,Q)_y^2 = (aQ)^2\{a_x Q_y + Q_x a_y\} = 0;
\]
anderntheils ist:
\[
(aQ)^2\{a_x Q_y - Q_x a_y\} = (aQ)^3(xy).
\]
Folglich erhält man durch Addition:
\[
(aQ)^2 a_x Q_y = \frac{1}{2}(aQ)^3(xy),
\]
und wenn man hierin $x$ durch $\Delta$ ersetzt, und $(aQ)^3$ durch seinen Werth $A_{\Delta\Delta}$, so kommt endlich:
\[
(a\Delta)(aQ)^2 \Delta_y Q_y = \frac{1}{2}A_{\Delta\Delta}\cdot \Delta_y^2;
\]
also ist in der That:
\weq{VI}{(Q,\,Q)^2 = \frac{1}{2}A_{\Delta\Delta}\cdot \Delta.}


% ===== pdf idx 186, printed 172.0 =====
% !PRINTED 172
\weberpage{172}
Weil aber auch aus bekannten Gründen $(Q,\,Q) = 0$ und $(Q,\,Q)^3 = 0$, so ist $(Q,\,Q')$ Reducent.

Demnach sind alle überhaupt möglichen Klammerfactoren Reducenten, und es besteht somit der Satz:

\begin{quote}
\glqq Das volle System der Form dritten Grades $f$ besteht aus den vier Formen $f$, $\Delta$, $A_{\Delta\Delta}$, $Q$.\grqq{}
\end{quote}

148. \textit{Tabelle der Ueberschiebungswerthe dieser vier Formen.} Ich will die in diesem Beweise gewonnenen Resultate noch einmal tabellarisch ordnen. Ist
\[
f = a_x^3,
\]
so hat man als symbolische Darstellungen der Grundformen des Systemes:
\[
(f,\,f)^2 = \Delta \phantom{{}={}} = (ab)^2 a_x b_x
\]
\[
(\Delta,\,\Delta)^2 = A_{\Delta\Delta} = R \phantom{{}={}} = (ab)^2(cd)^2(ad)(bc)
\]
\[
(f,\,\Delta) = Q \phantom{{}={}} = (c\Delta)c_x^2 \Delta_x = (ab)^2(cb)c_x^2 a_x.
\]
Als Ueberschiebungswerthe erhielten wir:
\weq{1}{(f,\,f) = 0, \quad (f,\,f)^2 = \Delta, \quad (f,\,f)^3 = 0}
\weq{2}{(\Delta,\,\Delta) = 0, \quad (\Delta,\,\Delta)^2 = A_{\Delta\Delta} = R}
\weq{3}{(Q,\,Q) = 0, \quad (Q,\,Q)^2 = \frac{1}{2}A_{\Delta\Delta}\cdot \Delta, \quad (Q,\,Q)^3 = 0}
\weq{4}{(f,\,\Delta) = Q, \quad (f,\,\Delta)^2 = 0}
\weq{5}{(\Delta,\,Q) = \frac{1}{2}A_{\Delta\Delta}\cdot f, \quad (\Delta,\,Q)^2 = 0}
\weq{6}{(f,\,Q) = -\frac{1}{2}\Delta^2, \quad (f,\,Q)^2 = 0, \quad (f,\,Q)^3 = A_{\Delta\Delta}.}
Unter den Covarianten des Systemes befindet sich eine schiefe Form, nämlich die Covariante $Q$. Ihr Quadrat muss sich durch die übrigen geraden Formen darstellen lassen. Um dies zu bewerkstelligen, berücksichtigen wir, dass sie Functionaldeterminante von $f$ und $\Delta$ ist. Indem wir die allgemeine Formel für Functionaldeterminanten benutzen (Nr. 48, I), gelangen wir zu der bereits von Cayley entdeckten Relation zwischen den vier Grundformen des Systemes:
\[
2Q^2 =
\begin{vmatrix}
(f,\,f)^2, & (f,\,\Delta)^2, & f \\
(f,\,\Delta)^2, & (\Delta,\,\Delta)^2, & \Delta \\
f, & \Delta, & 0
\end{vmatrix}
=
\begin{vmatrix}
\Delta, & 0, & f \\
0, & R, & \Delta \\
f, & \Delta, & 0
\end{vmatrix}
\]
oder:
\weq{I}{2Q^2 + \Delta^3 + Rf^2 = 0.}


% ===== pdf idx 187, printed 173.0 =====
% !PRINTED 173
\weberpage{173}
Wir werden bei Auflösung der Gleichung dritten Grades von dieser Relation Gebrauch machen.

149. \textit{Das Formensystem von $Q + \lambda f$.} Wir sahen, dass die Form $f$ eine Covariante gleichen Grades in den Variabeln besitzt, nämlich die Form $Q$ und können uns demnach die Frage vorlegen: welchen Werth haben die Grundformen des Büschels
\[
Q + \lambda f,
\]
also die Formen
\[
\Delta_{Q+\lambda f}, \quad R_{Q+\lambda f}, \quad Q_{Q+\lambda f}.
\]
Wir können zu diesem Zwecke den Weg directer Ueberschiebung oder auch den \gesperrt{Aronhold'schen} Process benutzen. Wir wollen hier den ersten einschlagen, und erhalten in einfacher Weise:
\weq{1}{\Delta_{Q+\lambda f} = (Q+\lambda f,\, Q+\lambda f)^2 = (Q,\,Q)^2 + 2\lambda(Q,\,f)^2 + \lambda^2(f,\,f)^2 = \frac{1}{2}R\cdot\Delta + \lambda^2\Delta = \left(\lambda^2 + \frac{R}{2}\right)\Delta}
\weq{2}{R_{Q+\lambda f} = \left(\left(\lambda^2+\frac{R}{2}\right)\Delta,\, \left(\lambda^2+\frac{R}{2}\right)\Delta\right)^2 = \left(\lambda^2+\frac{R}{2}\right)^2(\Delta,\,\Delta)^2 = \left(\lambda^2+\frac{R}{2}\right)^2 R}
\weq{3}{Q_{Q+\lambda f} = \left(Q+\lambda f,\, \left(\lambda^2+\frac{R}{2}\right)\Delta\right) = \left(\lambda^2+\frac{R}{2}\right)\left\{(Q,\,\Delta) + \lambda(f,\,\Delta)\right\} = \left(\lambda^2+\frac{R}{2}\right)\left\{-\frac{R}{2}f + \lambda Q\right\}.}

150. \textit{Beispiele zur Uebung.} Wir hatten zur Berechnung der Tabelle in Nr. 148 lediglich den Ueberschiebungsprocess benutzt. Der Umstand, dass $f$ eine Covariante $Q$ besitzt, die gleichfalls dritten Grades in den Variabeln ist, legt den Gedanken nahe, mit Hilfe des verwandten \gesperrt{Aronhold'schen} Processes solche Ueberschiebungen in vielleicht einfacherer Weise zu ermitteln. In der That können wir, indem wir eine beliebige Ueberschiebung von bekanntem Werthe deltairen, auch Werthe von andern Ueberschiebungen berechnen. Ich erinnere zunächst daran, dass wir in Nr. 63, unter der Voraussetzung
\weq{I}{\delta f = Q,}
folgende Beziehungen erhielten:
\weq{II}{\delta Q = -\frac{1}{2}Rf, \quad \delta\Delta = 0, \quad \delta R = 0.}
Deltairen wir nun die beiden Relationen
\[
\begin{aligned}
(f,\,\Delta)^2 &= 0\\
(f,\,Q)^2 &= 0,
\end{aligned}
\]
so erhalten wir:
\weq{1}{\delta(f,\,\Delta)^2 = (\delta f,\,\Delta)^2 + (f,\,\delta\Delta)^2 = 0,}
%CONT>>


% ===== pdf idx 188, printed 174.0 =====
% !PRINTED 174
\weberpage{174}
%<<CONT
also wegen den Beziehungen (I) und (II)
\[
(Q,\,\Delta)^2 = 0
\]
\weq{2}{\delta(f,\,Q)^2 = (\delta f,\,Q)^2 + (f,\,\delta Q)^2 = 0,}
oder:
\[
(Q,\,Q)^2 + \left(f,\,-\frac{1}{2}R\cdot f\right)^2 = 0,
\]
oder:
\[
(Q,\,Q)^2 = \frac{R}{2}\cdot\Delta.
\]
Man erkennt, wie auf diesem Wege oft in sehr einfacher Weise Ueberschiebungen berechnet werden können. Wir werden bei Gelegenheit der Theorie biquadratischer Formen noch ausgiebiger von dieser Methode Gebrauch machen.

151. \textit{Erste Anwendung: Die Discriminante der Form $f$.} Bekanntlich ist die Discriminante einer Form darstellbar als Resultante der gleich null gesetzten ersten Differentialquotienten dieser Form. (Vergl. Bd. I, Nr. 173 und 178.) In unserm Falle hier sind die beiden ersten Differentialquotienten von $f = a_x^3 = b_x^3$ zwei quadratische Formen
\[
f_1 = a_1 a_x^3
\]
\[
f_2 = b_2 b_x^2,
\]
% checked against the GDZ scan (27 Oct 2026): the print has f_1 = a_1 a_x^3 (exponent 3, a misprint: the text calls f_1, f_2 quadratic and has f_2 = b_2 b_x^2); reproduced as printed
und die Resultante derselben ist nach Nr. 129 gleich der Discriminante ihrer Functionaldeterminante $\vartheta$, nämlich:
\weq{1}{2A_{\vartheta\vartheta} = -R_{f_1f_2}.}
Nun ist die Functionaldeterminante von $f_1$ und $f_2$ dargestellt durch
\[
\vartheta = (ab)\,a_x b_x a_1 b_2,
\]
oder, wenn man $a$ mit $b$ vertauscht:
\[
\vartheta = -(ab)\,a_x b_x b_2 a_1,
\]
folglich:
\[
2\vartheta = (ab)^2 a_x b_x = \Delta,
\]
d. h.: \glqq Die Discriminante der Form $f$ ist also bis auf einen Zahlenfactor gleich der Discriminante von $\Delta$.\grqq{}
Man hat:
\[
A_{\vartheta\vartheta} = \frac{1}{4}A_{\Delta\Delta},
\]
also wegen Relation (1)
\[
-\frac{1}{2}A_{\Delta\Delta} = R_{f_1f_2}.
\]
Verschwindet die Discriminante von $f$, dann hat also sowohl $f$ als $\Delta$ zwei gleiche Wurzeln. Wir werden sofort sehen, dass in diesem Falle $\Delta$ ein rationaler Factor von $f$ ist, und dass die Form $Q$ den
%CONT>>


% ===== pdf idx 189, printed 175.0 =====
% !PRINTED 175
\weberpage{175}
%<<CONT
Doppelfactor der beiden Formen $f$ und $\Delta$ dreifach enthält, also ein reiner Cubus ist.

152. \textit{Zweite Anwendung: Auflösung der cubischen Gleichung.} Die Methoden, eine Gleichung dritten Grades aufzulösen, werden naturgemäss zunächst darauf gerichtet sein, sie auf eine reine cubische Gleichung, deren Wurzeln uns ja bekannt sind, zurückzuführen. Hiezu bieten uns die vorausgegangenen Entwicklungen zwei Relationen, einmal die \gesperrt{Cayley'sche} Identität
\weq{1}{2Q^2 + \Delta^3 + Rf^2 = 0,}
dann aber auch die Relation
\weq{2}{\Delta_{Q+\lambda f} = \left(\frac{R}{2}+\lambda^2\right)\cdot\Delta.}
Wir gehen zunächst von der zweiten Gleichung aus. Sobald nämlich die \gesperrt{Hesse'sche} Form $\Delta_{Q+\lambda f}$ von $Q+\lambda f$ identisch verschwindet, ist nach Nr. 55 die Form $Q+\lambda f$ ein reiner Cubus. Dies tritt also immer ein, sobald
\weq{3}{\frac{R}{2}+\lambda^2 = 0,}
und damit haben wir eine quadratische Resolvente der cubischen Form. Setzen wir ihre beiden Wurzeln:
\[
+\sqrt{-\frac{R}{2}} = \lambda_1
\]
\[
-\sqrt{-\frac{R}{2}} = \lambda_2,
\]
so sind die beiden Formen:
\[
Q+\lambda_1 f, \quad Q+\lambda_2 f
\]
die dritten Potenzen zweier linearer Formen, die wir mit $\alpha_x$ resp. $\beta_x$ bezeichnen wollen. Für jene Werthe von $x$ aber, welche das Verschwinden von $f$ bewirken, besteht wegen
\weq{4}{\left.\begin{aligned} Q+\lambda_1 f &= \alpha_x^3\\ Q+\lambda_2 f &= \beta_x^3 \end{aligned}\right\}}
die Beziehung
\[
\alpha_x^3 = \beta_x^3,
\]
d. h.: \glqq Die drei Wurzeln von $f$ bestimmen sich beziehungsweise aus den drei Gleichungen
\[
\begin{aligned}
\alpha_x - \beta_x &= 0\\
\alpha_x - \varepsilon\beta_x &= 0\\
\alpha_x - \varepsilon^2\beta_x &= 0.
\end{aligned}
\]\grqq{}
Demnach lässt sich $f$ bis auf eine Constante als Product folgender drei linearer Factoren schreiben:
%CONT>>


% ===== pdf idx 190, printed 176.0 =====
% !PRINTED 176
\weberpage{176}
%<<CONT
\weq{5}{\begin{aligned}
f = c\cdot\Bigg\{&\sqrt[3]{Q+\sqrt{-\frac{R}{2}}\cdot f} - \sqrt[3]{Q-\sqrt{-\frac{R}{2}}\cdot f}\Bigg\}\times\\
\Bigg\{&\sqrt[3]{Q+\sqrt{-\frac{R}{2}}\cdot f} - \varepsilon\sqrt[3]{Q-\sqrt{-\frac{R}{2}}\cdot f}\Bigg\}\times\\
\Bigg\{&\sqrt[3]{Q+\sqrt{-\frac{R}{2}}\cdot f} - \varepsilon^2\sqrt[3]{Q-\sqrt{-\frac{R}{2}}\cdot f}\Bigg\}.
\end{aligned}}
Es ist hiebei nicht nöthig, direct die dritte Wurzel aus $Q\pm\sqrt{-\frac{R}{2}}\cdot f$ auszuziehen; vielmehr braucht man nur die erste Polare zu bilden
\[
\begin{aligned}
Q_x^2 Q_y + \lambda_1 f_x^2 f_y &= \alpha_x^2\alpha_y\\
Q_x^2 Q_y + \lambda_2 f_x^2 f_y &= \beta_x^2\beta_y
\end{aligned}
\]
und nun etwa $x_1 = 1,\ x_2 = 0$ zu setzen. Die sich so ergebenden Werthe $\alpha_y$ und $\beta_y$ sind bis auf Constante gleich den dritten Wurzeln aus dem vollständigen Cubus $Q+\lambda_i f$.

153. Wir können nun aber auch die \gesperrt{Cayley'sche} Relation zur Auflösung der Gleichung dritten Grades benutzen, welche uns zugleich auch Aufschluss über die Bedeutung der beiden linearen Formen $\alpha_x$ und $\beta_x$ giebt, auf die wir im Vorausgehenden geführt wurden. Wir schreiben zu dem Zwecke die Relation (I) Nr. 148 in der Form:
\[
\Delta^3 = -\{2Q^2 + Rf^2\},
\]
oder auch
\[
\Delta^3 = -2\left\{Q+f\sqrt{-\frac{R}{2}}\right\}\left\{Q-f\sqrt{-\frac{R}{2}}\right\}.
\]
Auf der linken Seite befindet sich nun ein reiner Cubus, und rechts das Product zweier Formen dritten Grades, welche, da im allgemeinen $f$ und $Q$ keinen gemeinschaftlichen Factor besitzen werden, ebenfalls keinen Theiler haben. Demnach muss jeder Factor rechts einzeln ein voller Cubus sein, etwa:
\[
\begin{aligned}
\frac{1}{2}\left\{Q+f\sqrt{-\frac{R}{2}}\right\} &= \alpha_x^3\\
\frac{1}{2}\left\{Q-f\sqrt{-\frac{R}{2}}\right\} &= \beta_x^3,
\end{aligned}
\]
wo nun $\alpha_x$ und $\beta_x$ bis auf den Factor $\frac{1}{2}$ mit den in Nr. 152 gleich bezeichneten linearen Factoren übereinstimmen. Wir haben demnach
\[
\Delta^3 = -8\alpha_x^3\beta_x^3,
\]
oder
\weq{1}{\Delta = -2\alpha_x\beta_x.}


% ===== pdf idx 191, printed 177.0 =====
% !PRINTED 177
\weberpage{177}

Hiebei ist rechts noch eine der dritten Einheitswurzeln als Factor zu denken.

Wenn wir nun die Form $f$ so transformiren, dass an Stelle der Variabeln $x_1$ und $x_2$ die Linearfactoren von $\Delta$ eintreten, so geht $f=0$ in eine reine cubische Gleichung über. Es ist nämlich nach Nr. 148:
\[
(f,\Delta)^2=a_x(a\alpha)(a\beta)=0,
\]
d. h.: \glqq Jedes symbolische Product mit dem Factor $(a\alpha)(a\beta)$ hat den Werth null.\grqq{} Erhebt man daher die Identität
\[
a_x(\alpha\beta)=\alpha_x(a\beta)-\beta_x(a\alpha)
\]
auf die dritte Potenz, so ergiebt sich:
\weq{2}{a_x^3(\alpha\beta)^3=\alpha_x^3(a\beta)^3-\beta_x^3(a\alpha)^3=c\cdot f.}

Die Form $f$ ist also unter Zugrundelegung der Linearfactoren von $\Delta$ in eine reine cubische Form transformirt, und man erkennt zugleich, dass es hiebei wegen der überall auftretenden dritten Potenzen gleichgiltig ist, welche dritte Einheitswurzel wir der rechten Seite der Gleichung (1) als Factor geben.

Zur Auflösung der cubischen Gleichung:
\[
f=ax_1^{\,3}+3bx_1x_2^{\,2}+3cx_1x_2^{\,2}+dx_2^{\,3}=0
\] % [sic] both middle terms print the exponent 2 on x_2 as shown
hat man demnach folgenden Weg einzuschlagen.

Man bilde zunächst die Hesse'sche Form von $f$
\[
\frac{1}{2}\Delta=(ac-b^2)x_1^{\,2}+(ad-bc)x_1x_2+(bd-c^2)x_2^{\,2}
\]
und suche die beiden Wurzeln derselben:
\[
\begin{matrix}\lambda_i\\ \mu_i\end{matrix}
=\frac{(bc-ad)\pm\sqrt{(ad-bc)^2-4(bd-c^2)(ac-b^2)}}{2(ac-b^2)}.
\]

Es ist alsdann:
\[
\frac{1}{2}\Delta=\frac{1}{\mu_1}(x_1\mu_1-x_2\lambda_1)(x_1\mu_2-x_2\lambda_2).
\]
Mit Hilfe der Transformation
\[
\begin{aligned}
y_1&=\mu_1x_1-\lambda_1x_2=\alpha_x\\
y_2&=\mu_2x_1-\lambda_2x_2=\beta_x
\end{aligned}
\]
geht $f$ in die Form über
\[
F=Ay_1^{\,3}+Dy_2^{\,3}=0,
\]
die nun durch Ausziehen der dritten Wurzel gelöst werden kann.

154. \textit{Schlussbemerkungen.} Sind die Coefficienten von $f$ reell und ist $R=A_{\Delta\Delta}$ negativ, so besitzt $f$ eine reelle und zwei imaginäre Wurzeln. Dies erkennt man direct aus der Darstellung (5) der Form $f$ in Nr. 152.


% ===== pdf idx 192, printed 178.0 =====
% !PRINTED 178
\weberpage{178}

Ist $R=0$, so besitzt $f$ zwei gleiche Wurzeln und ebenso $\Delta$. In diesem Falle ist natürlich die in Nr. 153 gegebene Transformation von $f$ unmöglich. % checked against the GDZ scan (27 Oct 2026): the raised mark after "unmöglich." is an irregular ink blob, not a footnote sign (the footnote anchors of this book are "*)", and the page has no footnote); not reproduced
Die Hesse'sche Form ist dann, wie schon früher erwähnt, Factor von $f$. Denn da in jedem Falle nach Nr. 152 (4)
\[
\begin{aligned}
Q+\sqrt{-\frac{R}{2}}\,f&=\alpha_x^3\\
Q-\sqrt{-\frac{R}{2}}\,f&=\beta_x^3,
\end{aligned}
\]
so ist demnach:
\[
\begin{aligned}
2f\cdot\sqrt{-\frac{R}{2}}&=\alpha_x^3-\beta_x^3=(\alpha_x-\beta_x)\cdot(\alpha_x-\varepsilon\beta_x)\cdot(\alpha_x-\varepsilon^2\beta_x)\\
2Q&=\alpha_x^3+\beta_x^3\\
\Delta&=\varrho\cdot\alpha_x\beta_x,
\end{aligned}
\]
wo $\varrho$ eine Constante.

Ist nun $R=0$, also $\alpha_x=\beta_x$, so wird
\[
\begin{aligned}
Q&=\alpha_x^3\\
\Delta&=\varrho\alpha_x^2.
\end{aligned}
\]
Es ist aber
\[
(f,\Delta)=Q,
\]
oder:
\[
\varrho(a\alpha)a_x^2\alpha_x=\alpha_x^3,
\]
also nach Division mit dem linearen Factor $\alpha_x$:
\[
\varrho(a\alpha)a_x^2=\alpha_x^2.
\]
Daraus schliesst man, dass $\alpha_x$ der Doppelfactor von $f$. [Vergl. Nr. 99, (3).] Ueberdies sieht man, dass, im Falle $R=0$, $Q$ die dritte Potenz des Doppelfactors von $f$ ist.

Ist endlich $\Delta$ selbst identisch null, so ist nach Nr. 55 die Form $f$ der Cubus einer linearen Form.

\begin{center}
\S\ 15. \textbf{Das Formensystem der Form vierten Grades.}
\end{center}

155. \textit{Die fundamentalen In- und Covarianten.} Wie bei den Formen dritten Grades, so liefern auch hier die Ueberschiebungen von $f$ über sich selbst und über ihre Hesse'sche Covariante das vollständige System der Form vierten Grades, d. h. die Gesammtheit jener Formen, durch welche sich alle übrigen Co- und Invarianten ganz und rational ausdrücken lassen. Wir haben diese Ueberschiebungen $(f,f)^\lambda$, $(f,\Delta)^\lambda$ bereits früher, insbesondere auch in § 8 unter den Anwendungen des Processes der Reihenentwicklung, kennen gelernt, und wollen nur noch einmal die dort gewonnenen Resultate zusammenfassen.


% ===== pdf idx 193, printed 179.0 =====
% !PRINTED 179
\weberpage{179}

Bezeichnen wir mit
\[
\begin{aligned}
f&=\bar a_0x_1^{\,4}+4\bar a_1x_1^{\,3}x_2+6\bar a_2x_1^{\,2}x_2^{\,2}+4\bar a_3x_1x_2^{\,3}+\bar a_4x_2^{\,4}\\
&=a_x^4=b_x^4=c_x^4=\dots\ \text{etc.}
\end{aligned}
\]
die gegebene Form vierten Grades, so sind die erwähnten Formen durch folgende symbolische Producte dargestellt:
\[
\begin{aligned}
(f,f)^2&=(ab)^2a_x^2b_x^2=\Delta\\
(f,f)^4&=(ab)^4\phantom{a_x^2b_x^2}=i\ {}^{\ast})\\
(f,\Delta)&=(a\Delta)a_x^3\Delta_x^3=(ab)^2(cb)c_x^3a_x^2b_x=t\\
(f,\Delta)^4&=(a\Delta)^4\phantom{a_x^2b_x^2}=(ab)^2(bc)^2(ca)^2=j\ {}^{\ast}).
\end{aligned}
\]
Diese fünf Formen:
\[
f,\ \Delta,\ t,\ i,\ j
\]
bilden, wie schon erwähnt, das vollständige System der Form $f$, und mit dem Beweise dieser Behauptung wollen wir uns nun im Folgenden beschäftigen.

156. \textit{Gedankengang des Beweises.} Irgend eine Co- oder Invariante von $f$ kann nur durch Ueberschiebung der Formen $f$, $\Delta$, $t$ über sich selbst oder über einander entstehen. Legen wir hiebei $t$ in der Form $(a\Delta)a_x^3\Delta_x^3$ zu Grunde, so enthält jedes symbolische Product, das bei diesem Ueberschiebungsprocess entsteht, nur Symbole $a$ und $\Delta$, und zwar ausser in Factoren erster Art $a_x$, $\Delta_x$ nur in den drei allein möglichen Klammerfactoren:
\[
(ab),\ (a\Delta),\ (\Delta\Delta_1).
\]
Gelingt es uns also zu zeigen, dass alle Formen, welche einen dieser drei Klammerfactoren enthalten, reducibel sind, vergl. Nr. 124, so ist damit auch bewiesen, dass ausser den erwähnten fünf Formen keine weitern existiren können, die sich nicht rational und ganz durch die gegebenen $f$, $\Delta$, $t$, $i$, $j$ ausdrücken lassen.

Wir werden demnach zuerst die zunächst liegenden einfachen Ueberschiebungen von $f$ über $\Delta$, und $\Delta$ über $\Delta$ berechnen, wobei

\par\noindent\rule{4cm}{0.4pt}\par
\noindent $^{\ast})$ Es ist:
\[
i=(ab)^4=a_1^4b_2^4-4a_1^3a_2\cdot b_2^3b_1+6a_1^2a_2^2b_1^2b_2^2-\cdots\ \text{etc.}=2(\bar a_0\bar a_4-4\bar a_1\bar a_3+3\bar a_2^2),
\]
\[
j=(ab)^2(bc)^2(ca)^2=\begin{vmatrix}a_1^2 & a_1a_2 & a_2^2\\ b_1^2 & b_1b_2 & b_2^2\\ c_1^2 & c_1c_2 & c_2^2\end{vmatrix}^2=6\begin{vmatrix}\bar a_0 & \bar a_1 & \bar a_2\\ \bar a_1 & \bar a_2 & \bar a_3\\ \bar a_2 & \bar a_3 & \bar a_4\end{vmatrix}
\]
\[
=6(\bar a_0\bar a_2\bar a_4-\bar a_0\bar a_3^2-\bar a_1^2\bar a_4-\bar a_2^3+2\bar a_1\bar a_2\bar a_3)\ \text{(vergl. Bd. I, Nr. 40).}
\]


% ===== pdf idx 194, printed 180.0 =====
% !PRINTED 180
\weberpage{180}
%<<CONT
sich als erstes Resultat ergiebt, dass sich die so entstehenden Formen in der That durch die fünf Grundformen rational und ganz ausdrücken lassen.

Sodann zeigen wir auf Grund dieser Thatsache, dass die Klammerfactoren
\[
(ab),\ (a\Delta)^2,\ (a\Delta)(b\Delta),\ (\Delta\Delta_1)
\]
Reducenten sind, und daraus ergiebt sich alsdann, dass jedes durch Ueberschiebung von $f$, $\Delta$ und $t$ entstehende Product reducibel ist.

Die obenerwähnten Klammerfactoren $(ab)$ und $(\Delta\Delta_1)$ sind unter diesen Reducenten, und was symbolische Producte mit dem Factor $(a\Delta)$ betrifft, so haben sie entweder diesen ganz allein und dann nothwendig auch noch die Factoren $a_x^3\cdot\Delta_x^3$, d. h. sie besitzen $t$ als Factor; oder sie haben als nächst einfachsten Fall die Factoren $(a\Delta)^2$, resp. $(a\Delta)(b\Delta)$, und von diesen wollen wir ja gerade beweisen, dass sie Reducenteneigenschaft besitzen.

157. \textit{Die Ueberschiebungen $(f,\Delta)^\lambda$, $(\Delta,\Delta)^\lambda$.} Den ersten Theil des Beweises haben wir bereits in früheren Untersuchungen zum grössten Theil erledigt, oder können ihn wenigstens auf Grund der dort gewonnenen Resultate rasch zum Ziele führen.

Die Ueberschiebungen $(f,\Delta)$, $(f,\Delta)^4$ haben wir in das System aufgenommen, und $(\Delta,\Delta)$, $(\Delta,\Delta)^3$ verschwinden identisch; die Formen $(f,\Delta)^2$, $(f,\Delta)^3$ wurden in § 8 berechnet als erste Anwendung des Processes der Reihenentwicklung, wobei wir die Werthe
\[
\begin{aligned}
(f,\Delta)^2&=\frac{i}{6}\cdot f\\
(f,\Delta)^3&=0
\end{aligned}
\]
erhielten. Es erübrigt also nur noch die Berechnung von
\[
(\Delta,\Delta)^2\ \text{und}\ (\Delta,\Delta)^4.
\]

158. \textit{Berechnung von $(\Delta,\Delta)^2$.} Die Ueberschiebung $(\Delta,\Delta)^2$ geht aus der zweiten Polare $\Delta_{y^2}$ hervor, wenn wir darin $y$ durch $\Delta_1$ ersetzen, und mit $\Delta_{1x}^2$ multipliciren. Diese zweite Polare ist aber das erste Glied in der Reihenentwicklung der Form
\[
(ab)a_x^3b_y^3=\frac{3}{2}\Delta_{y^2}(xy)+\frac{1}{4}i(xy)^3.
\]
Daher ist:
\weq{1}{(ab)(b\Delta)^3a_x^3\Delta_x=\frac{3}{2}(\Delta,\Delta)^2+\frac{1}{4}i\cdot\Delta.}

Hiebei entsteht das symbolische Product links, wenn wir in $(b\Delta)^3b_y\Delta_x$, $y$ durch $a$ ersetzen und mit $a_x^3$ multipliciren. Entwickeln
%CONT>>


% ===== pdf idx 195, printed 181.0 =====
